\documentclass[12pt]{article}

% Match the earlier reports using locally available fonts.
\usepackage[T1]{fontenc}
\usepackage{times}
\usepackage{microtype}
\usepackage{amsmath}
\usepackage[parfill]{parskip}
\setlength{\parindent}{0pt}
\setlength{\parskip}{2ex plus 0.2ex minus 0.2ex}
\usepackage[letterpaper,margin=1in]{geometry}
\usepackage[hidelinks]{hyperref}

\begin{document}

\begin{flushleft}
\textbf{Week 7: Matrix Factorization and Efficient MAP Inference} \\
Jaisel Singh (js6897) \\
October 28, 2026
\end{flushleft}

\vspace{0.1in}

What stood out to me in Koren, Bell, and Volinsky's review was how a simple
shared representation can predict a large collection of ratings. In the basic
model, a user and an item each have a short vector, and their dot product
predicts the rating. The connection to maximum a posteriori (MAP) inference
becomes clear if I view squared error as a Gaussian likelihood and quadratic
regularization as Gaussian priors on the factors. We obtain a point estimate
of those factors, rather than the distribution approximated in last week's
variational inference. Alternating least squares also makes the computation
concrete: fixing one set of factors turns fitting the other into quadratic
problems, although the joint objective remains nonconvex.

Lee and Seung made me think more carefully about what the factors mean.
Nonnegative factors reconstruct a face by adding components, so features cannot
cancel each other through subtraction. Their examples make the role of
constraints in producing interpretable representations especially clear.
Kolda and Bader add a useful distinction: some higher-order tensor
decompositions are unique up to scaling and permutation under conditions that
do not hold for ordinary matrix factorizations. I find it helpful to separate
that mathematical identifiability from semantic interpretation. Even a unique
decomposition would not establish that a component corresponds to a meaningful
human preference.

Gopalan, Hofman, and Blei show how the observation model also shapes computation.
Their hierarchical Poisson factorization models consumption with nonnegative
user and item factors, while accounting for differences in user activity and
item popularity. I found the treatment of sparse data particularly useful:
zero entries still contribute to the likelihood, but their contribution can be
aggregated through sums of factors without visiting every user--item pair.
Their fitting procedure uses variational inference, connecting this application
directly to last week. It makes scalability feel like a consequence of modeling
choices as well as implementation. My remaining question is how to check
whether factors that predict consumption well capture preferences, or mainly
reflect which items users had an opportunity to encounter.

\vspace{0.05in}
{\footnotesize
\textit{Readings:} Koren et al.\ (2009), \href{https://www.cs.columbia.edu/~blei/fogm/readings/KorenBellVolinsky2009.pdf}{\emph{Matrix Factorization Techniques for Recommender Systems}};
Lee and Seung (1999), \href{https://www.cs.cmu.edu/~efros/courses/LBMV07/Papers/lee-nature-99.pdf}{\emph{Learning the Parts of Objects by Non-negative Matrix Factorization}};
Gopalan et al.\ (2015), \href{https://www.cs.columbia.edu/~blei/papers/GopalanHofmanBlei2015.pdf}{\emph{Scalable Recommendation with Hierarchical Poisson Factorization}};
Kolda and Bader (2009), \href{https://www.cs.columbia.edu/~blei/fogm/readings/KoldaBader2009.pdf}{\emph{Tensor Decompositions and Applications}}.
\par}

\end{document}

%%% Local Variables:
%%% mode: latex
%%% TeX-master: t
%%% End:
